\documentclass[border=10pt]{standalone}
\usepackage{fontspec}
\AtBeginDocument{\hyphenation{zenodo}}% never hyphenate "zenodo" inside a DOI
\setmainfont{Noto Serif}[Ligatures=TeXOff]
\XeTeXgenerateactualtext=1
\newfontfamily\bengalifont{Noto Serif Bengali}[Script=Bengali, NFSSFamily=bnfam]
\usepackage[Latin,Bengali]{ucharclasses}
% group-free font switching (as in header-common.tex): robust inside braces, cells and quotes
\makeatletter
\def\BN@saved{\rmdefault}
\def\BN@bn{bnfam}
\newcommand\BNin{\edef\BN@tmp{\f@family}\ifx\BN@tmp\BN@bn\else\let\BN@saved\BN@tmp\fi\fontfamily{bnfam}\selectfont}
\newcommand\BNout{\fontfamily{\BN@saved}\selectfont}
\makeatother
\setTransitionsFor{Bengali}{\BNin}{\BNout}
% Typographic quotes and dashes are not in a ucharclasses class (class 0), so a Bangla word right
% after “ or ‘ would not trigger the Bangla font. Treat them as Latin so that every switch fires.
\XeTeXcharclass"2018=\XeTeXcharclass`A
\XeTeXcharclass"2019=\XeTeXcharclass`A
\XeTeXcharclass"201C=\XeTeXcharclass`A
\XeTeXcharclass"201D=\XeTeXcharclass`A
\XeTeXcharclass"2013=\XeTeXcharclass`A
\XeTeXcharclass"2014=\XeTeXcharclass`A
\XeTeXcharclass"2026=\XeTeXcharclass`A
\usepackage{newunicodechar}
\newunicodechar{।}{{\bengalifont\char"0964}}
\newunicodechar{॥}{{\bengalifont\char"0965}}

\usepackage[table]{xcolor}
\usepackage{booktabs,tabularx,array}
\definecolor{bapilrule}{HTML}{1F2937}
\definecolor{bapilstripe}{HTML}{F3F4F6}
\renewcommand{\arraystretch}{1.28}
\setlength{\tabcolsep}{5pt}

\begin{document}
\begin{minipage}{17cm}
{\large\bfseries সারণি ১৯।}\ {\large রেফারেন্স স্থাপত্যের উপাদান ও সম্ভাব্য প্রযুক্তি}\par\vspace{7pt}
\small
\rowcolors{2}{bapilstripe}{white}
\begin{tabularx}{\linewidth}{@{}>{\hsize=0.309\hsize\raggedright\arraybackslash}X>{\hsize=0.979\hsize\raggedright\arraybackslash}X>{\hsize=1.546\hsize\raggedright\arraybackslash}X>{\hsize=1.340\hsize\raggedright\arraybackslash}X>{\hsize=0.825\hsize\raggedright\arraybackslash}X@{}}
\arrayrulecolor{bapilrule}\toprule
\rowcolor{white}\textbf{আইডি} & \textbf{উপাদান} & \textbf{কাজ} & \textbf{সম্ভাব্য উন্মুক্ত প্রযুক্তি} & \textbf{রেফারেন্স বাস্তবায়ন} \\
\midrule
C1 & উৎস সংযোগকারী & সাইটের শর্ত মেনে নির্ধারিত সময়ে জনসমক্ষে প্রকাশিত নথি সংগ্রহ; উৎস-তথ্য সংরক্ষণ & Python worker, Scrapy, cron বা Airflow & পরিকল্পিত \\
C2 & গ্রহণ ও প্রমিতকরণ & PDF ও HTML থেকে টেক্সট নিষ্কাশন, বাংলা OCR, ইউনিকোড প্রমিতকরণ, খণ্ডকরণ & Apache Tika [65], Tesseract (ben, eng) [66], নিজস্ব প্রমিতকারী & bapil.text, bapil.corpus \\
C3 & নথি সংরক্ষণাগার ও ক্যাটালগ & মূল নথি, নিষ্কাশিত টেক্সট ও DCAT-সামঞ্জস্যপূর্ণ মেটাডেটা, সংস্করণসহ & PostgreSQL [67], MinIO (S3-সামঞ্জস্যপূর্ণ), CKAN [68] & চেকসামসহ JSONL কর্পাস \\
C4 & অনুসন্ধান সূচি & শব্দভিত্তিক BM25 ও ঘন ভেক্টরের মিশ্র র‍্যাংকিং & OpenSearch [69], বা pgvector-সহ [70] PostgreSQL & bapil.search (BM25, দ্বিভাষিক) \\
C5 & উত্তর সেবা & উদ্ধার, উদ্ধৃতি, বিরত থাকা; উদ্ধৃতি যাচাইসহ ঐচ্ছিক LLM উৎপাদন & vLLM বা llama.cpp-এ চালিত উন্মুক্ত-ওজন LLM & bapil.answer (উদ্ধৃতিসহ নিষ্কাশনমূলক) \\
C6 & বিশ্লেষণ ইঞ্জিন & বাজেট নিষ্কাশন, সূচক, বিষয়-মডেল & Python, pandas, DuckDB & পরিকল্পিত \\
C7 & পোর্টাল ও ড্যাশবোর্ড & দ্বিভাষিক ও প্রবেশযোগ্য ব্যবহারকারী ইন্টারফেস & ওয়েব ফ্রেমওয়ার্ক; Apache Superset & পরিকল্পিত \\
C8 & API গেটওয়ে & সংস্করণকৃত REST API, কোটা, হার-সীমা & FastAPI, Kong & bapil.api (স্ট্যান্ডার্ড লাইব্রেরি) \\
C9 & পরিচয় ও প্রবেশাধিকার & একক সাইন-অন, ভূমিকাভিত্তিক প্রবেশাধিকার, বহু-উপাদান প্রমাণীকরণ & Keycloak & পরিকল্পিত \\
C10 & নিরীক্ষা ও পর্যবেক্ষণযোগ্যতা & হস্তক্ষেপ-শনাক্তযোগ্য লগ, মেট্রিক, ট্রেসিং & OpenTelemetry, Prometheus, Grafana & bapil.audit (হ্যাশ-শৃঙ্খলিত JSONL) \\
C11 & মডেল রেজিস্ট্রি ও মূল্যায়ন & মডেল কার্ড, সংস্করণকৃত মূল্যায়ন, বেঞ্চমার্ক & MLflow & bapil.evaluate \\
C12 & শাসন কর্মক্ষেত্র & এআই ব্যবস্থা নিবন্ধন, ঝুঁকি মূল্যায়ন, ঘটনা, স্বচ্ছতা নিবন্ধন & হালকা ওয়েব অ্যাপ্লিকেশন & bapil.risk, bapil.registry \\
\bottomrule
\end{tabularx}
\par\vspace{5pt}{\footnotesize\itshape প্রযুক্তিগুলো সম্ভাব্য বিকল্প, বাধ্যতামূলক নয়; চূড়ান্ত পছন্দ কারিগরি স্থাপত্য বোর্ড NFR-02, NFR-03 ও NFR-06 অনুযায়ী নির্ধারণ করবে।}
\par\vspace{3pt}{\scriptsize\color{gray} উৎস: BAPIL টেকনিক্যাল ব্লুপ্রিন্ট, সংস্করণ ১.০ (২০২৬)। CC BY 4.0।}
\end{minipage}
\end{document}
